%! Unit 5: Determinants and Linear Transformations — Summative Assessment — Practice Test --- Study Copy (blank)
\documentclass[10pt]{article}
\usepackage{linalg-article}
\usepackage{linalg-boxes}
\newcommand{\xx}{\mathbf{x}}
\newcommand{\ee}{\mathbf{e}}
\newcommand{\T}{^{\mathsf T}}

% Part-divider strip
\newcommand{\parthead}[1]{%
  \vspace{6pt}\noindent
  \begin{headlinebox}{burgundy}{\color{white}\bfseries #1}\end{headlinebox}%
  \vspace{4pt}}

\begin{document}

\pageheader{Unit 5: Determinants and Linear Transformations}{Practice Test --- Study Copy}
\namedateperiod

\vspace{0.10in}
\begin{remindbox}
\small
\textbf{This is a practice test.} It mirrors the real Unit 5 test in length, format,
and the ideas it covers, but the numbers and contexts differ. Work every problem
with no notes, then check your answers against the key.
\end{remindbox}

% ======================================================================
\parthead{Part A --- Vocabulary (8 pts)}
Write the letter of the best description next to each term.

\vspace{4pt}
\noindent
\begin{minipage}[t]{0.40\textwidth}
\begin{enumerate}[leftmargin=2.2em, itemsep=7pt, label=\arabic*.]
  \item Determinant \blank{1.2cm}
  \item Minor \blank{1.2cm}
  \item Cofactor expansion \blank{1.2cm}
  \item Linear transformation \blank{1.2cm}
  \item Orientation \blank{1.2cm}
  \item Singular matrix \blank{1.2cm}
  \item Area/volume factor \blank{1.2cm}
  \item Product rule \blank{1.2cm}
\end{enumerate}
\end{minipage}%
\hfill
\begin{minipage}[t]{0.57\textwidth}
\small
\begin{description}[leftmargin=1.4em, itemsep=3pt]
  \item[(A)] The smaller determinant left after crossing out one entry's row and column.
  \item[(B)] A map $\xx\mapsto A\xx$ that fixes the origin and sends grid lines to straight, evenly spaced lines.
  \item[(C)] A matrix with $\det=0$: its columns are dependent, it collapses space, and it has no inverse.
  \item[(D)] A number computed from a square matrix; its absolute value is the area (or volume) of the box built from the columns.
  \item[(E)] The rule $\det(AB)=\det A\,\det B$: composing transformations multiplies their scaling factors.
  \item[(F)] Computing a determinant along one row: each entry times its signed minor, with the pattern $+\,-\,+\,\dots$
  \item[(G)] Whether a transformation keeps or flips the handedness of space --- kept when $\det>0$, flipped when $\det<0$.
  \item[(H)] The quantity $|\det A|$, the number every area (or volume) is multiplied by under the transformation $A$.
\end{description}
\end{minipage}

% ======================================================================
\parthead{Part B --- Multiple Choice (12 pts)}
Circle the best answer.

\begin{enumerate}[leftmargin=1.9em, itemsep=9pt]
  \item The determinant of the identity matrix $I$ is
    \hfill (A) $0$ \quad (B) $1$ \quad (C) $n$ \quad (D) undefined
  \item Swapping two rows of a matrix
    \hfill (A) leaves $\det$ unchanged \quad (B) negates $\det$ \quad (C) doubles $\det$ \quad (D) makes $\det=0$
  \item Adding a multiple of one row to another row
    \hfill (A) negates $\det$ \quad (B) scales $\det$ \quad (C) leaves $\det$ unchanged \quad (D) makes $\det=0$
  \item A square matrix with $\det=0$ is
    \hfill (A) invertible \quad (B) singular \quad (C) orthogonal \quad (D) symmetric
  \item The number $|\det A|$ tells you
    \hfill (A) the trace \quad (B) how much $A$ scales area/volume \quad (C) the rank \quad (D) nothing geometric
  \item The determinant of a product satisfies
    \hfill (A) $\det(AB)=\det A+\det B$ \quad (B) $\det(AB)=\det A\,\det B$ \quad (C) $\det(AB)=\det A-\det B$ \quad (D) no rule
\end{enumerate}

% ======================================================================
\parthead{Part C --- Short Answer \& Computation (35 pts)}
Show your work.

\begin{enumerate}[leftmargin=1.9em, itemsep=12pt]
  \item Compute the determinant of $A=\begin{bsmallmatrix}1&3\\2&1\end{bsmallmatrix}$. Then give the
        \textbf{area} of the parallelogram built from its columns, and say whether $A$
        \textbf{preserves or reverses orientation}.
        \vspace{1.8cm}
  \item Compute the determinant of $B=\begin{bsmallmatrix}1&2&0\\0&3&1\\1&0&2\end{bsmallmatrix}$ by
        \textbf{cofactor expansion} across the top row (show the three $2\times2$ minors).
        \vspace{2.6cm}
  \item Compute the determinant of $C=\begin{bsmallmatrix}1&2&1\\2&5&3\\0&2&5\end{bsmallmatrix}$ by
        \textbf{reducing to triangular form} (elimination), then multiplying the pivots.
        \vspace{2.8cm}
  \item A $3\times3$ matrix $A$ has $\det A=6$. Find each new determinant and name the row rule you used:
        \textbf{(a)} two rows of $A$ are swapped; \textbf{(b)} one row of $A$ is multiplied by $3$;
        \textbf{(c)} $2\times(\text{row }1)$ is added to row $3$; \textbf{(d)} $\det(2A)$.
        \vspace{2.6cm}
  \item Two $3\times3$ matrices have $\det A=4$ and $\det B=2$. Use the product rule to find
        \textbf{(a)} $\det(AB)$, \textbf{(b)} $\det(A^{-1})$, \textbf{(c)} $\det(A\T)$, \textbf{(d)} $\det(A^2)$.
        \vspace{2.4cm}
  \item Photo software applies the transformation $\xx\mapsto A\xx$ with $A=\begin{bsmallmatrix}3&1\\0&2\end{bsmallmatrix}$.
        \textbf{(a)} What is the area of the image of the unit square? \textbf{(b)} A triangle has area $4$;
        what is the area of its image? \textbf{(c)} Does the image keep its orientation?
        \vspace{2.4cm}
  \item Let $M=\begin{bsmallmatrix}2&1\\4&2\end{bsmallmatrix}$. \textbf{(a)} Compute $\det M$.
        \textbf{(b)} Is $M$ invertible? \textbf{(c)} Describe what the transformation $\xx\mapsto M\xx$ does
        to the unit square (what does it collapse to, and why).
        \vspace{2.4cm}
\end{enumerate}

% ======================================================================
\parthead{Part D --- Extended Response (10 pts)}
Explain your reasoning in full sentences.

\begin{enumerate}[leftmargin=1.9em, itemsep=16pt]
  \item You compute the determinant of a $3\times3$ matrix, then perform one elimination step ---
        adding a multiple of one row to another --- and compute the determinant again.
        \textbf{(a)} Explain why the determinant does \emph{not} change. \textbf{(b)} Using the picture of the
        box (parallelepiped) built from the rows, explain why sliding the box sideways this way keeps its
        volume the same.
        \vspace{3.0cm}
  \item A transformation $\xx\mapsto A\xx$ sends the unit square to a parallelogram whose sides are the
        columns $A\ee_1$ and $A\ee_2$. \textbf{(a)} Explain why the area of that parallelogram is $|\det A|$,
        and why this makes $|\det A|$ the factor by which \emph{every} area is scaled. \textbf{(b)} Explain
        why $\det A=0$ forces $A$ to have no inverse, using what happens to the unit square.
        \vspace{2.8cm}
\end{enumerate}

\end{document}
